\label{sec:real-line-one-dimensional}

\begin{topicbox}{The Line as Order Plus Distance}
\begin{exposition}
The real numbers carry two compatible structures at once: a linear order $<$ and a notion of distance $d(x,y)=|x-y|$. Together these make $\mathbb{R}$ a one-dimensional space --- a line. This chapter develops only the metric topology of $\mathbb{R}$, not topology in general spaces. Every open set is built from intervals, and compactness is introduced only on $\mathbb{R}$, with Heine--Borel as the capstone.
\end{exposition}
\begin{exposition}
The intuition is geometric: points are positions on a line, order is ``left of,'' and distance is the length of the segment between two points. The caution to keep throughout is that $\mathbb{R}$ is special. Order and metric agree here in a way that fails on a general space, so definitions are stated in the form that survives to metric spaces and manifolds, with $\mathbb{R}$ as the first instance rather than the only one.
\end{exposition}
\end{topicbox}
